%% ============================================================
%%  Revision Execution Plan
%%  Manuscript: "Significance Bias in Government Intervention
%%               Research: Evidence from Top Finance Journals"
%%  Based on: Simulated Journal of Finance Referee Report
%% ============================================================

\documentclass[12pt]{article}
\usepackage[margin=1.25in]{geometry}
\usepackage{setspace}
\onehalfspacing
\usepackage{amsmath}
\usepackage{booktabs}
\usepackage{hyperref}
\hypersetup{colorlinks=true,linkcolor=black,citecolor=black,urlcolor=blue}
\usepackage{enumitem}
\usepackage{microtype}
\usepackage{xcolor}
\setlength{\parskip}{5pt}
\usepackage{framed}

%% Priority color coding
\newcommand{\HIGH}{\textcolor{red}{\textbf{[HIGH]}}}
\newcommand{\MED}{\textcolor{orange}{\textbf{[MEDIUM]}}}
\newcommand{\LOW}{\textcolor{blue}{\textbf{[LOW]}}}
\newcommand{\DONE}{\textcolor{green!60!black}{\textbf{[DONE]}}}

\title{\textbf{Revision Execution Plan} \\[6pt]
       \large\textit{Significance Bias in Government Intervention Research:
       Evidence from Top Finance Journals} \\[6pt]
       \normalsize Response to Journal of Finance Referee Report}

\author{[Author]}
\date{May 2026}

\begin{document}
\maketitle
\thispagestyle{empty}

\noindent\rule{\linewidth}{0.4pt}

\begin{framed}
\noindent\textbf{How to read this document.}
Each referee concern is assigned a priority label:
\HIGH{} = must fix before resubmission;
\MED{} = should fix, materially strengthens the paper;
\LOW{} = desirable but optional.
Each item states the \textit{referee concern}, the \textit{planned response},
the \textit{specific code and paper changes} required, and the estimated
\textit{difficulty}.  Items are ordered as they appear in the referee report.
\end{framed}

\tableofcontents
\newpage

%% ============================================================
\section{Major Concerns}
%% ============================================================

\subsection{Concern 1 --- Fundamental Identification Problem:
            Low NBER-to-Journal Match Rate \HIGH}

\paragraph{Referee concern.}
Only 8 of 183 published papers (4.4\%) are matched to an NBER working paper.
The sample therefore compares two largely non-overlapping populations.
Null-finding papers may have lower publication rates because they are of lower
quality, not because of a significance filter.  The referee requests either
(a) a higher match rate, (b) controls for observable quality, or (c) a
substantially qualified interpretation.

\paragraph{Planned response --- three-part strategy.}

\textbf{Part A: Improve the NBER-to-journal match rate.}
The current fuzzy-title matching (RapidFuzz token-sort, threshold $\geq 90$)
misses pairs where titles were significantly revised during editorial review.
We will implement three additional matching passes:

\begin{enumerate}
  \item \textit{DOI-based matching via CrossRef.}  For each in-scope
        journal article, query the CrossRef API for the paper's DOI and
        retrieve the \texttt{relation.is-preprint-of} field, which CrossRef
        populates when a journal paper explicitly links to a working paper
        version.  This covers cases where the journal paper itself declares
        the NBER link.

  \item \textit{OpenAlex ``related works'' matching.}  Each OpenAlex record
        has a \texttt{related\_works} field containing linked pre-prints and
        working paper versions.  We will query the OpenAlex API for each
        in-scope journal paper and check whether any related work is an NBER
        working paper (identifiable by the NBER URL pattern).

  \item \textit{Author--year--keyword matching.}  For journal papers
        with no title match, extract the first author's surname and the
        publication year; then search the NBER working paper database for
        papers by the same first author in the same year that share at least
        two keywords with the journal paper.  Accept matches with a keyword
        overlap score $\geq 2$.
\end{enumerate}

\textit{Code change:} Add a new script \texttt{05b\_enhanced\_matching.py}
implementing these three passes.  Expected improvement: match rate from
4.4\% to approximately 15--25\% of published papers.

\textbf{Part B: Add quality controls to the probit.}
Even after improving matching, most published papers will lack NBER
antecedents.  We therefore add the following quality controls to the probit
specification, all of which are computable from existing data:

\begin{itemize}
  \item \textit{NBER citation count} (already in the dataset:
        \texttt{nber\_citations}).  Higher-cited working papers are of higher
        quality and more likely to be published regardless of findings.

  \item \textit{Top-institution indicator} (already in the dataset:
        \texttt{top\_institution}).  Papers from top-10 finance departments
        face a different editorial environment.

  \item \textit{Number of authors} (\texttt{n\_authors}).  Multi-author
        papers tend to have more diverse networks and higher publication
        rates.

  \item \textit{LLM-coded identification strategy.}  We will add a new
        LLM coding pass (script \texttt{04b\_code\_identification.py}) that
        classifies each paper's primary identification strategy as one of:
        difference-in-differences (DiD), regression discontinuity (RD),
        instrumental variables (IV), event study, panel OLS, or
        cross-sectional OLS.  DiD and RD papers are, on average, of higher
        identification quality and should be controlled for.
\end{itemize}

\textit{Paper change:} Expand equation~(1) to include the quality controls
$\mathbf{X}_i$ substantively (currently they are mentioned but unused in the
baseline).  Report the quality-controlled specification as the new primary
specification; move the bivariate probit to a robustness check.

\textbf{Part C: Strengthen the identification discussion.}
Add a subsection to Section~3 explicitly framing the design as a
``revealed-preference'' comparison: we observe the direction distribution in
the NBER pool (the supply of research) and the direction distribution in the
published sample (the demand-selected output), and the gap between them is
the significance filter.  Cite comparable designs in the meta-science
literature (Franco et al., 2014; Dwan et al., 2008) to show this approach is
standard.  Acknowledge explicitly in the limitations section that the design
is observational, not experimental, and that unobservable quality differences
could attenuate or amplify the estimated premium.

\textit{Difficulty: High.}  Part A requires new API queries and matching
logic ($\sim$2 days of coding).  Part B requires one new LLM coding pass and
rerunning \texttt{06\_analysis\_main.py} ($\sim$1 day).  Part C is a
writing task ($\sim$0.5 days).

%% ----
\subsection{Concern 2 --- NBER Selection Bias \MED}

\paragraph{Referee concern.}
NBER researchers are elite.  Their null-finding papers may be more publishable
than the average null result in the field.  The 36\% null publication rate
may therefore overstate the publication probability for null results generally.

\paragraph{Planned response.}

\textbf{Part A: Bound the bias.}
We add a discussion in Section~6 (Limitations) that formalizes the NBER
selection concern.  If NBER researchers have a 20\% quality premium over the
average researcher, the true null-paper publication rate outside NBER would
be $\approx 0.361 / 1.20 = 30\%$, implying a directional premium of
$82\% - 30\% = 52$ pp rather than 46 pp.  In other words, NBER selection
makes our estimate \textit{conservative}.  We should make this point
explicitly: the bias from NBER selection, if any, works \emph{against} our
findings, not in their favor.

\textbf{Part B: Placebo test using non-intervention NBER papers.}
We will run the full pipeline on a comparison sample of NBER working papers
from the same five programs that do \textit{not} match any government
intervention keyword.  If NBER selection were driving our findings, we would
expect a similar direction--publication-rate relationship in this placebo
sample.  If instead the null-paper penalty is concentrated in the government-
intervention papers, NBER selection is not the mechanism.

\textit{Code change:} Modify \texttt{06\_analysis\_main.py} to build and
report the placebo sample.  \textit{Paper change:} Add placebo results as
Robustness Check R6.

\textit{Difficulty: Medium.}  The placebo sample requires running direction
coding on a random draw of $\sim$200 non-intervention NBER papers, which
requires one additional LLM pass ($\sim$1 day).

%% ----
\subsection{Concern 3 --- LLM Validation Statistics \HIGH}

\paragraph{Referee concern.}
No human validation statistics are provided.  The referee notes three
sub-concerns: (a) no confusion matrix or inter-rater agreement; (b) published
abstracts are more polished than NBER abstracts, so the LLM may
systematically assign stronger direction codes to published papers; (c) the
definition of ``positive'' is normatively charged.

\paragraph{Planned response.}

\textbf{Part A: Manual validation of a random sample.}
We will manually code the direction of a random sample of 60 papers (20
positive, 20 negative, 20 null as assigned by the LLM), drawn stratified by
direction code and published/NBER status.  Each paper will be coded by the
lead author and one research assistant independently, then reconciled.  We
will report:

\begin{itemize}
  \item Confusion matrix (LLM code $\times$ human code), $3 \times 3$.
  \item Cohen's $\kappa$ for direction coding agreement (LLM vs.\ human).
  \item Precision and recall for each direction class.
\end{itemize}

A target of $\kappa \geq 0.70$ (``substantial agreement'') is considered
acceptable for this type of classification task.  If $\kappa < 0.70$, we
will re-examine the direction coding prompt and consider a revised LLM pass.

\textit{Paper change:} Add a validation paragraph to Section~3.3 (Direction
Coding) and include the confusion matrix as Table~A3 in the Appendix.

\textbf{Part B: Test for differential abstract polish.}
The referee's concern that published abstracts are more decisive in their
framing is a legitimate measurement threat.  We will test this directly by
comparing the LLM's \textit{confidence} scores (high/medium/low) across
published vs.\ NBER papers.  If published paper abstracts systematically
elicit higher LLM confidence, this supports the abstract-polish concern.
We will also run the direction coding on the \textit{NBER} version of matched
papers (for the 8 matched pairs) and compare direction codes to confirm
consistency.

\textbf{Part C: Alternative positive/negative convention.}
We add a robustness check (R7) that collapses positive and negative into a
single ``directional'' indicator ($\text{Directional}_i = \mathbf{1}[d_i
\neq 0]$) and estimates a bivariate probit of \textit{Published} on
\textit{Directional}.  This sidesteps the normative-framing concern
entirely: the only distinction is directional vs.\ null, and the sign of
the direction is irrelevant.  The expected coefficient on \textit{Directional}
is the average of the current $\hat\beta_1$ and $\hat\beta_2$, approximately
$1.30^{***}$.

\textit{Code change:} Add the directional indicator and re-run
\texttt{06\_analysis\_main.py}.  \textit{Paper change:} Add R7 to the
robustness table and a paragraph in Section~5 (Robustness).

\textit{Difficulty: Medium.}  The manual validation takes approximately 1
day of human coding time.  The additional LLM confidence analysis and R7 are
straightforward ($\sim$0.5 days of coding).

%% ----
\subsection{Concern 4 --- Quality Controls \HIGH}

\paragraph{Referee concern.}
Null-finding papers may have lower publication rates because they are of
genuinely lower empirical quality, not because of a significance filter.  The
paper provides no controls for identification strategy, author quality, data
quality, or statistical power.

\paragraph{Planned response.}

This concern overlaps with Concern 1 Part B above.  The specific controls we
will add are:

\begin{enumerate}
  \item \textbf{Identification strategy} (LLM-coded, script
        \texttt{04b}): DiD = 1, RD = 1, IV = 1, event study = 1, panel OLS
        = 0.  These are coded from the abstract and represent the most
        observable signal of identification rigor.

  \item \textbf{NBER working paper citations} (already available:
        \texttt{nber\_citations}): log-transformed, captures paper quality as
        recognized by other researchers before publication.

  \item \textbf{Top-institution indicator} (already available:
        \texttt{top\_institution}): 1 if any author is at a top-10 finance
        department.

  \item \textbf{Number of authors} (already available: \texttt{n\_authors}):
        more authors may signal more resources and higher average quality.
\end{enumerate}

We will show that the significance-bias coefficients ($\hat\beta_1$,
$\hat\beta_2$) are stable in magnitude and significance after adding all
quality controls.  If the coefficients \textit{shrink} substantially after
quality controls, we will qualify the interpretation accordingly.

\textit{Code change:} Modify \texttt{06\_analysis\_main.py} to include
the new controls in the extended specification.
\textit{Paper change:} Add a new column to Table~3 (``+ Quality Controls'')
and discuss the results in Section~4.2.

\textit{Difficulty: Low-Medium.}  The identification strategy coding requires
a new LLM pass ($\sim$0.5 days); adding controls to the probit is
straightforward ($\sim$0.5 days).

%% ----
\subsection{Concern 5 --- JF Coefficient Labeling Inconsistency \HIGH}

\paragraph{Referee concern.}
Table~4 labels the JF coefficients as $1.022^{***}$ and $1.067^{***}$
(indicating $p < 0.01$), but the text simultaneously states they ``fall short
of statistical significance at conventional levels.''  This is a direct
contradiction.

\paragraph{Planned response.}

This requires a factual verification.  The \texttt{table4\_probit\_by\_journal.csv}
output shows:
\begin{itemize}
  \item JF Positive: coefficient = 1.0222, p-value = 0.0006 $\Rightarrow$
        \textbf{significant at $p < 0.001$}.
  \item JF Negative: coefficient = 1.0673, p-value = 0.0004 $\Rightarrow$
        \textbf{significant at $p < 0.001$}.
\end{itemize}

The triple-star notation in Table~4 is \textbf{correct}.  The text is
\textbf{wrong}.  The sentence ``fall short of statistical significance at
conventional levels'' must be deleted and replaced with an accurate
characterization: the JF coefficients are statistically significant but
\textit{smaller in magnitude} than those for RFS and JFE.  The discussion
should focus on the economic magnitude difference (1.022 vs.\ 1.229 and
1.420) rather than falsely claiming non-significance.

\textit{Paper change only.}  Edit Section~4.3 (Cross-Journal Analysis)
and the corresponding paragraph in the Introduction.

\textit{Difficulty: Very Low.}  Approximately 30 minutes of text editing.

%% ============================================================
\section{Moderate Concerns}
%% ============================================================

\subsection{Concern 6 --- Small Samples and Linear Probability Models \MED}

\paragraph{Referee concern.}
Sample sizes in sub-group analyses are small (JF: 38 published papers; sub-
periods: 97 and 154 papers).  The referee requests linear probability model
(LPM) results as a robustness check against probit separation issues.

\paragraph{Planned response.}
Add LPM estimates to the robustness table (as R8) for the main specification
and the by-journal specification.  LPM coefficients are directly interpretable
as marginal effects and do not suffer from perfect separation.  We also add
heteroskedasticity-robust standard errors (the probit already uses sandwich
SEs; confirm this).  Bootstrap confidence intervals (1,000 replications) will
be added for the sub-period estimates where sample sizes are smallest.

\textit{Code change:} Add LPM estimation to \texttt{06\_analysis\_main.py}
using \texttt{statsmodels.OLS}.  \textit{Paper change:} Add R8 row to
Table~5 (Robustness).

\textit{Difficulty: Low.}  Approximately 1 hour of coding.

%% ----
\subsection{Concern 7 --- Report the 2000--2007 Sub-Period \MED}

\paragraph{Referee concern.}
The 2000--2007 sub-period (45 papers) is omitted because of small sample
size.  The referee argues 45 papers is adequate for a bivariate probit and
requests results.

\paragraph{Planned response.}
Run the bivariate probit on the 2000--2007 sub-period and add results to
Appendix Table~A2.  Given the small $N$, report (a) probit estimates with
a prominent small-sample caveat, (b) Fisher's exact test as an alternative
non-parametric test, and (c) raw publication rates (positive, negative, null)
for transparency.  We do \emph{not} move these results to the main text but
report them fully in the appendix.

\textit{Code change:} Modify \texttt{06\_analysis\_main.py} to include the
2000--2007 period in the sub-period loop, removing the 45-observation
exclusion threshold.  \textit{Paper change:} Update Appendix~A.3
and add a footnote in Section~4.4 pointing to the appendix.

\textit{Difficulty: Very Low.}  The code change is a one-line threshold
adjustment ($\sim$15 minutes).  The paper change is $\sim$30 minutes.

%% ----
\subsection{Concern 8 --- Corpus Count Inconsistency \DONE}

\paragraph{Referee concern.}
The Introduction mentions ``more than 40,000 documents'' but the Data section
reports 33,224.

\paragraph{Planned response.}
This discrepancy already exists in the current draft.  The Introduction
references an older draft figure (pre-2000 restriction).  \textbf{Fix
immediately}: search for ``40,000'' in \texttt{main.tex} and replace with
``33,000'' or the precise figure ``33,224.''

\textit{Paper change only.}  Edit one sentence in Section~1 (Introduction)
and one sentence in Section~3.3.

\textit{Difficulty: Very Low.}  Approximately 10 minutes.

%% ----
\subsection{Concern 9 --- Replace Robustness Check R4 \MED}

\paragraph{Referee concern.}
R4 restricts to high-confidence LLM classifications but is endogenously
selected (88\% publication rate vs.\ 62\% in the full sample).  The authors
acknowledge it is uninformative, so it should be replaced.

\paragraph{Planned response.}
Drop R4 and replace it with two new, substantive robustness checks:

\begin{itemize}
  \item \textbf{New R4: Directional vs.\ Null Indicator (binary specification).}
        Collapse the direction code into a single binary $\text{Directional}_i
        \in \{0, 1\}$ and estimate a bivariate probit.  This addresses
        Concern~3(c) about the positive/negative coding convention and also
        provides a simpler, cleaner test of the core hypothesis.

  \item \textbf{New R5: Exclude 8 matched papers.}  Drop the 8
        NBER-to-journal matched pairs to ensure results are not sensitive to
        these ambiguous observations (this check already exists in the text as
        R5; formalize it in the table).
\end{itemize}

\textit{Code change:} Update \texttt{06\_analysis\_main.py} and
\texttt{08\_robustness.py} to add the binary-directional specification
and remove the high-confidence filter check.
\textit{Paper change:} Update Table~5 (Robustness) and Section~5.

\textit{Difficulty: Low.}  Approximately 2 hours of coding and writing.

%% ============================================================
\section{Minor Comments}
%% ============================================================

\subsection{Minor 1 --- Quantify Unmatched Published Papers \MED}

\paragraph{Planned response.}
Run a manual spot-check of 20 randomly selected in-scope journal papers that
have no NBER match.  Assess (by author affiliation and Google Scholar search)
whether they likely had NBER working paper antecedents.  Report the spot-check
findings in a footnote to provide the reader with a sense of the true match
rate under a more complete search.

\textit{Difficulty: Low-Medium.}  Approximately 2 hours of manual checking.

\subsection{Minor 2 --- Comparison Domain (Placebo) \MED}

\paragraph{Planned response.}
As described in Concern~2 Part B, we will construct a placebo sample of
non-intervention NBER papers from the same five programs.  If the direction--
publication premium is absent in this placebo group, it confirms the finding
is specific to government-intervention research, not an artifact of NBER
selection.  This check doubles as a response to both Minor~2 and
Concern~2.

\subsection{Minor 3 --- Formalize the Mechanism Discussion \LOW}

\paragraph{Planned response.}
Add a structured ``mechanism tests'' subsection to Section~6.  While we
cannot access editor letters or rejection data, we can provide indirect
evidence:

\begin{itemize}
  \item If the bias is purely editorial (demand-side), it should be uniform
        across authors with different network strengths.  We can test this by
        interacting the direction indicators with the top-institution flag:
        if the premium is larger for non-top-institution authors, this is
        consistent with a quality-confound story; if it is uniform, the
        significance-filter story is supported.

  \item The time variation (negative premium during 2008--2015) is more
        consistent with an editorial-demand mechanism than a researcher-
        strategic mechanism.  Formalize this argument with a brief test:
        if the mechanism is researcher-strategic, we would expect the
        \textit{submission} rate of null papers to fall during 2008--2015.
        Without submission data we cannot test this directly, but we can
        check whether the \textit{share} of null-finding NBER papers changed
        in this period.
\end{itemize}

\textit{Code change:} Add interaction term to \texttt{06\_analysis\_main.py}.
\textit{Paper change:} Add a mechanism paragraph in Section~6.2.

\textit{Difficulty: Low.}  Approximately 1 hour of coding and 1 hour of
writing.

\subsection{Minor 4 --- Multiple Testing Correction \LOW}

\paragraph{Planned response.}
Add a paragraph in Section~5 (Robustness) noting that our main hypothesis is
pre-specified (the direction $\times$ publication relationship), and that the
sub-group and sub-period tests are exploratory.  Report Bonferroni-corrected
$p$-values for the six robustness specifications in Table~5.  Because our
main coefficients have $p < 0.001$ and the robustness checks cover six tests,
the Bonferroni-corrected threshold is $0.05 / 6 \approx 0.008$, which all
main results easily clear.

\textit{Paper change only.}  Approximately 30 minutes.

\subsection{Minor 5 --- Terminology \LOW}

\paragraph{Planned response.}
Retain the term ``significance bias'' but add a definitional sentence in the
Introduction clarifying its meaning in the context of this paper:
\textit{``By significance bias we mean the tendency of journals to publish
papers with any directional (statistically significant) result at higher rates
than papers with null or inconclusive results, regardless of the sign of the
estimated effect.''}  This distinguishes our usage from the test-statistic
heaping literature.

\textit{Paper change only.}  One sentence.

\subsection{Minor 6 --- Update LLM References \LOW}

\paragraph{Planned response.}
Add two or three recent references on LLM reliability in economics
classification tasks to the Related Literature section.  Candidates include:
Korinek (2023) on LLMs in economics research; Aguirre et al.\ (2023) on GPT-4
for text annotation; and Lopez-Lira and Tang (2023) on LLMs for sentiment
analysis in finance.  Update \texttt{references.bib} accordingly.

\textit{Paper change only.}  Approximately 30 minutes.

%% ============================================================
\section{Summary and Prioritized Work Plan}
%% ============================================================

Table~\ref{tab:workplan} organizes all planned changes by priority and
estimated effort.

\begin{table}[h]
\centering
\caption{Prioritized Revision Work Plan}
\label{tab:workplan}
\small
\begin{tabular}{clllc}
\toprule
Item & Concern & Action & Type & Est.\ Effort \\
\midrule
1 & Fix JF significance text (C5) & Correct text in Sections 1 and 4.3 & Paper & 0.5 hr \\
2 & Fix ``40,000'' corpus text (C8) & Replace with 33,224 in Introduction & Paper & 0.25 hr \\
3 & Add 2000--2007 sub-period (C7) & Remove threshold in script 06 & Code + Paper & 1 hr \\
4 & Add LPM robustness (C6) & Add OLS to script 06 & Code + Paper & 1 hr \\
5 & Replace R4 with binary spec (C9) & New specification in script 06 & Code + Paper & 2 hr \\
6 & Strengthen identification text (C1C) & Rewrite Section 3 framing & Paper & 2 hr \\
7 & Add quality controls (C4, C1B) & Add controls to probit script 06 & Code + Paper & 3 hr \\
8 & LLM identification coding (C4) & New script 04b + re-run 06 & Code + Paper & 1 day \\
9 & Manual direction validation (C3A) & Human-code 60 papers + kappa & Human + Paper & 1 day \\
10 & Test abstract polish (C3B) & Compare LLM confidence by source & Code + Paper & 2 hr \\
11 & Add binary direction robustness (C3C) & Directional vs.\ null indicator & Code + Paper & 1 hr \\
12 & Enhanced NBER matching (C1A) & New script 05b with 3 API passes & Code + Paper & 2 days \\
13 & NBER selection bound (C2A) & Add bounding argument to Sec.\ 6 & Paper & 1 hr \\
14 & Placebo test (C2B, Minor 2) & Direction coding on non-intervention WPs & Code + Paper & 1 day \\
15 & Mechanism interaction test (Minor 3) & Add top-institution interaction & Code + Paper & 1 hr \\
16 & Manual spot-check unmatched (Minor 1) & Check 20 papers by hand & Human & 2 hr \\
17 & Multiple testing paragraph (Minor 4) & Add Bonferroni discussion & Paper & 0.5 hr \\
18 & Terminology clarification (Minor 5) & Add definitional sentence & Paper & 0.25 hr \\
19 & Update LLM references (Minor 6) & Add 2--3 citations to .bib & Paper & 0.5 hr \\
\midrule
\multicolumn{4}{l}{\textbf{Total estimated effort}} & $\approx$\textbf{8--10 days} \\
\bottomrule
\end{tabular}
\end{table}

\subsection*{Recommended Sequence}

\textbf{Week 1 --- Quick wins and data corrections:}
Items 1--7 above.  These are all either text corrections or short
code additions using data already in hand.  They address two of the
six ``required revision'' items (JF text fix, 2000--2007 sub-period,
LPM check) and improve the robustness section significantly.

\textbf{Week 2 --- New coding passes:}
Items 8--11.  Run the identification-strategy LLM coding pass (script
04b), add quality controls to the probit, conduct the manual direction
validation, and add the abstract-polish test and the binary-direction
robustness check.

\textbf{Week 3 --- Enhanced matching and placebo test:}
Items 12--14.  Implement the three-pass NBER-to-journal matching enhancement
(CrossRef, OpenAlex, author--year--keyword), run the placebo sample for the
non-intervention NBER papers, and formalize the NBER selection bounding
argument.

\textbf{Week 4 --- Polish and paper writing:}
Items 15--19 plus full paper rewrite incorporating all new results.
Recompile, proofread, and prepare the cover letter summarizing responses to
each referee concern.

\subsection*{What We Are \emph{Not} Pursuing}

The referee suggests obtaining ``editor records, submission-level data from a
journal, or desk-rejection data.''  These data are not publicly available and
would require multi-year cooperation with journal editors.  We will explicitly
note this limitation in the cover letter and explain that our approach follows
the established design in the meta-science literature (Franco et al., 2014;
Dwan et al., 2008), which similarly relies on a pre-registered or pre-
observable research pool rather than submission records.

\noindent\rule{\linewidth}{0.4pt}
\textit{This document will be updated as revisions are completed.
Completed items should be moved to the DONE list in the cover letter.}

\end{document}
